\section*{Chapitre \ref{ch:iv3:system-design} --- Conception de systèmes}
\markboth{\protect\omTocMark{}{Réponses~: Conception de systèmes}}{\protect\omTocMark{}{Réponses~: Conception de systèmes}}

\begin{solution}{iq:iv3:system-design:1}
En moyenne, 2 sont dans le système, puisque $50\,000 \times 40 \times 10^{-6} = 2$. À une profondeur de 200, le temps moyen est $200 / 50\,000 = 4$ ms. C'est la loi de Little, $L = \lambda W$~: le nombre moyen d'éléments dans le système est égal au taux d'arrivée multiplié par le temps moyen passé dans le système.

Elle suppose seulement un système stable sur une période assez longue, de sorte que ce qui entre finisse par sortir. Elle ne demande aucune hypothèse sur la loi des arrivées ni sur celle des temps de service, ce qui explique son utilité. Le test la vérifie sur une file simulée avec des arrivées aléatoires.

En pratique, elle permet de mesurer deux grandeurs et d'en déduire la troisième. La profondeur de file se lit à peu de frais en production, alors que le temps passé dans le système demande plus de précautions. Une profondeur de 200 à un débit connu indique aussitôt que les messages ont 4~ms lorsqu'ils sortent.

\iqwrong{«~Les 40 microsecondes sont un coût fixe, donc une profondeur de 200 ne change rien.~» Le temps passé dans le système inclut l'attente, et c'est elle qui a augmenté.}
\iqfollow{Cela s'applique-t-il pendant une rafale~?}{Seulement en moyenne sur une fenêtre assez longue. Pendant une rafale, il entre davantage qu'il ne sort, si bien que l'hypothèse de système stable n'est plus vérifiée à cet instant.}
\end{solution}

\begin{solution}{iq:iv3:system-design:2}
Pour 10 backends, $1 - 0.99^{10} = 9.6~\%$. Pour 100 backends, $1 - 0.99^{100} = 63~\%$. Une requête qui attend la réponse de tous les backends est aussi lente que le plus lent d'entre eux, si bien qu'un événement rare sur chacun devient fréquent au sommet de l'arbre d'appels. La moyenne masque cela, car une faible part d'appels lents la déplace à peine.

Le centile d'un seul backend est donc ce que ressentent ses utilisateurs, et cela se répète à chaque appel. Une étude publiée sur de grands services web donne le même chiffre de 63~\% pour ce calcul, à 100 serveurs.

Les remèdes viennent du fait de ne pas attendre tout le monde. Envoyer une seconde copie d'une requête lente vers un autre réplica après un court délai, et retenir la première réponse. Renvoyer un résultat partiel si la conception le permet. Réduire le fan-out, ou supprimer les causes des requêtes lentes, soit 1~\%, comme les pauses du ramasse-miettes et les tâches de fond.

L'indépendance est une hypothèse. Si la lenteur a une cause commune, comme un incident réseau, moins de requêtes sont touchées, puisque les appels lents arrivent groupés, mais sur celles-ci tous les backends sont lents en même temps.

\iqwrong{«~Chaque backend respecte sa cible de 99~\% à 5~ms, donc le service aussi.~» Les cibles ne se composent pas à travers un fan-out.}
\end{solution}

\begin{solution}{iq:iv3:system-design:3}
Le pic demande $10^6 \times 40 \times 8 = 320$ Mb/s. Une journée représente $200\,000 \times 40 \times 6.5 \times 3600 = 187$ Go, et une année de 252 jours environ 47~To. Il faut dimensionner tout ce qui traite des données en direct sur le pic, et le stockage sur la moyenne.

Un chemin de traitement temps réel qui prend du retard au pic est en échec, quelle que soit sa moyenne. Les liens, les handlers et les consommateurs sont donc dimensionnés pour le pic, avec une marge, car les pics croissent d'année en année et surviennent aux pires moments. Le stockage et les traitements par lots de nuit s'intéressent au total de la journée, donc ils utilisent la moyenne.

Sur quelle fenêtre le pic est-il mesuré~? Un pic à la seconde masque des rafales bien plus marquées à l'échelle de la milliseconde, qui est celle que voit un tampon. Il faut donc demander le pic à la plus fine échelle disponible. Sur une milliseconde, un débit peut valoir plusieurs fois le chiffre à la seconde.

Ce sont des chiffres bruts. Les en-têtes s'ajoutent au débit de ligne, et la compression réduit le stockage, d'un facteur qu'il faut mesurer sur ses propres données.

\iqwrong{«~La moyenne est de 200\,000 par seconde, donc 64 Mb/s suffisent.~» Une liaison à ce débit perdrait des paquets ou les mettrait en file pendant chaque minute chargée.}
\end{solution}

\begin{solution}{iq:iv3:system-design:4}
L'indisponibilité vaut $0.001 \times 8760 = 8.76$ heures par an. Si les deux sont nécessaires, on obtient $0.999^2 = 99.8~\%$, soit environ 17.5 heures d'arrêt. Si l'un suffit, on obtient $1 - 0.001^2 = 99.9999~\%$, soit environ 31.5 secondes d'arrêt. L'indépendance est le point faible, car les pannes réelles partagent des causes.

Les deux règles sont la multiplication des disponibilités en série, et celle des indisponibilités en parallèle. Ainsi, chaque dépendance ajoutée dans une chaîne dégrade la situation, et la redondance l'améliore beaucoup, sur le papier.

Deux copies sont rarement indépendantes. Elles partagent une alimentation électrique, un commutateur, une release portant le même bogue, une poussée de configuration, ou un certificat qui expire le même jour. Une seule mauvaise release les met donc toutes deux à l'arrêt en même temps, et les 31 secondes par an relèvent de la fiction. La bascule (failover) prend aussi du temps, et la logique de bascule peut elle-même tomber en panne, généralement au moment où on en a besoin pour la première fois.

Une revue de conception doit donc chercher les causes communes plutôt que compter les neufs. Il faut utiliser des domaines de panne distincts, déployer les changements d'un côté à la fois et tester la bascule.

\iqwrong{«~Deux réplicas donnent six neufs.~» Ce n'est que le calcul sous hypothèse d'indépendance. Un bogue commun ne donne qu'un seul neuf.}
\end{solution}

\begin{solution}{iq:iv3:system-design:5}
Il affiche \texttt{2 [100, 100]} et \texttt{1 [100]}. Sans clé, le nouvel envoi est un second ordre, et l'on se retrouve avec deux fois la position voulue. Avec un identifiant d'ordre client, la passerelle reconnaît le nouvel envoi et renvoie la première réponse. Le nouvel envoi devient alors sans effet de bord, c'est ce qu'on appelle l'idempotence.

Un délai d'attente n'apprend rien au client sur l'ordre~: celui-ci peut ne jamais être arrivé, ou seule la réponse peut s'être perdue. Un client doit donc réémettre ou interroger, et une clé est ce qui rend cela sûr.

La passerelle doit conserver une clé pendant au moins toute la durée où un client pourrait la réémettre. En trading, cela correspond en général à une session ou à une journée, et les identifiants sont uniques à l'intérieur de celle-ci.

L'ordre et la réponse mémorisée doivent former une seule écriture atomique. Si la passerelle tombe après avoir accepté l'ordre et avant d'avoir enregistré la clé, un nouvel envoi crée un second ordre. Si elle enregistre d'abord la clé, un nouvel envoi reçoit une réponse pour un ordre qui n'a jamais existé. Les deux doivent donc aller dans la même entrée de journal, ou dans la même transaction.

\iqwrong{«~Réémettre au dépassement du délai, et TCP garantit l'absence de doublons.~» TCP élimine les doublons au sein d'une même connexion. Un nouvel envoi après reconnexion est un nouveau message.}
\end{solution}

\begin{solution}{iq:iv3:system-design:6}
Sur le chemin critique (hot path)~: la décision, le contrôle de risque, l'envoi de l'ordre et la mise à jour de la position, puisque la décision suivante et le contrôle de risque ont besoin de la position. Hors du chemin critique~: l'enregistrement d'audit, la ligne de journal et le tableau de bord. Ceux-ci ne doivent pas pouvoir retarder l'ordre suivant.

Les autres composants reçoivent leurs données par une file, en général un tampon circulaire (ring buffer)~: le thread du chemin critique écrit un petit enregistrement binaire de taille fixe, et un autre thread ou processus fait le reste. Même le formatage d'un horodatage ou d'un double est laissé à ce consommateur. Ainsi, le chemin critique ne coûte qu'une copie en mémoire et une écriture.

Cela a une conséquence qu'il faut bien avoir en tête. Si la file se remplit, il faut bien céder quelque chose~: faut-il abandonner des enregistrements de journal, ou arrêter de trader~? Un enregistrement d'audit exigé par le régulateur ou pour son propre rapprochement ne doit pas être abandonné~; une file pleine à cet endroit est donc une raison de cesser d'envoyer des ordres, et de déclencher une alerte. Un flux vers le tableau de bord peut simplement être abandonné.

Le contrôle de risque reste sur le chemin critique, puisqu'il doit pouvoir arrêter l'ordre. Il doit rester simple et s'exécuter en temps constant.

\iqwrong{«~Tout journaliser en ligne, puisque la journalisation coûte peu.~» La moyenne est faible, mais elle peut se bloquer sur un disque ou un verrou au pire moment, et c'est cette queue de distribution qu'on observera.}
\end{solution}

\begin{solution}{iq:iv3:system-design:7}
Le polling coûte au serveur 10\,000 réponses par seconde dans les deux cas. Le push en coûte 20\,000 par seconde si le prix change deux fois par seconde, et environ 167 par seconde s'il change une fois par minute. Le retard moyen du polling vaut la moitié de l'intervalle, soit 0.5~s. Le polling est le meilleur choix lorsque les changements sont plus fréquents que ce dont les clients ont besoin.

Le polling fixe le rythme selon ce que veut le client, et le push le fixe selon ce que font les données. Ainsi, avec des données calmes, le push coûte bien moins cher et fournit des valeurs plus fraîches. Avec des données actives, il peut saturer le client, à moins de recourir à la conflation. Et un client en polling qui interroge une fois par seconde voit au pire un prix vieux de 1~s, si bien que sa charge et sa fraîcheur sont faciles à analyser.

Le polling présente aussi des atouts pratiques. Il ne demande aucun état par client sur le serveur ni aucune connexion ouverte, et il se rétablit seul après n'importe quelle panne, puisque l'interrogation suivante obtient la valeur courante. Un système en push doit gérer les reconnexions et les mises à jour manquées.

Un compromis consiste à pousser un instantané conflaté à cadence fixe, ce que demande souvent vraiment un affichage de prix.

\iqwrong{«~Le push est toujours plus efficace.~» Avec un prix qui change deux fois par seconde, il a envoyé deux fois plus de messages ici.}
\end{solution}

\begin{solution}{iq:iv3:system-design:8}
Le maximum est de 500\,000 ordres par seconde. À 400\,000, la charge $\rho$ vaut 0.8, et le temps moyen dans le système est de $2 / (1 - 0.8) = 10$ microsecondes, dont 8 en attente. À une charge de 100~\%, la file n'a pas de régime stationnaire, si bien que les attentes croissent sans limite. L'exécution du test à 110~\% de charge le montre.

Cette formule vaut pour des arrivées aléatoires et un service exponentiel, elle n'est donc qu'un modèle d'un thread réel. Elle en montre toutefois la forme~: la latence dépend de $1 / (1 - \rho)$, si bien que les derniers pourcents de capacité coûtent le plus cher.

Les arrivées sont aussi par rafales. Un débit de 400\,000 moyenné sur une seconde peut contenir des instants bien au-dessus de 500\,000. Une file se forme alors quoi qu'il arrive, et il faut une capacité libre pour la vider. Un thread qui fonctionne en moyenne près de son maximum ne se vide jamais.

Il faut donc dimensionner la capacité de sorte que le pic laisse une vraie marge, et surveiller le comportement de la file, pas seulement le débit moyen. Le débit maximal est le point où le système s'effondre, et non un régime de fonctionnement.

\iqwrong{«~Il tient 500\,000, donc 490\,000 est sûr.~» Sous ce modèle, l'attente moyenne y serait d'environ cinquante fois le temps de service.}
\end{solution}

\begin{solution}{iq:iv3:system-design:9}
Avec la règle du modulo, environ 10/11, soit 91~\%, des clés changent de nœud, car presque chaque reste change. Avec un anneau, environ 1/11, soit 9~\%, bougent, et toutes vont sur le nouveau nœud. Sur notre exécution de 100\,000 clés, les deux méthodes donnent 90.7~\% et 9.8~\%. Le \Cref{lst:iv3:sd:ring} présente l'anneau.
\omcode{code/interviews-3/06-system-design/python/iv_ring.py}{1}{22}{Un anneau de hachage.}\label{lst:iv3:sd:ring}

Chaque nœud est haché sur un cercle, et une clé appartient au point suivant dans le sens horaire. Un nouveau nœud ne reprend que les arcs situés juste avant ses propres points, si bien que les autres clés restent en place. Le test vérifie que chaque clé déplacée atterrit sur le nouveau nœud.

Avec un seul point par nœud, les arcs ont des longueurs très inégales. Dans notre exécution, le nœud le plus chargé parmi les 10 portait 1.94 fois la moyenne. Avec 100 nœuds virtuels, ce rapport tombe à 1.11, et avec 1000 à 1.04. Les nœuds virtuels répartissent aussi la charge d'un nœud défaillant sur tous les autres, et non sur un seul voisin.

Le sharding par hachage convient mal lorsque la charge est très inégale. Les symboles ne se valent pas~: quelques-uns ont beaucoup plus d'activité que les autres, et un hachage ne le voit pas, si bien qu'un même serveur peut recevoir plusieurs symboles très actifs. En trading, on place souvent les symboles délibérément en fonction de la charge mesurée, et l'on regroupe les instruments liés, de sorte que les deux jambes d'un spread vivent sur un même thread. La stabilité de l'affectation pendant une session compte davantage que la commodité.

\iqwrong{«~Le hachage répartit uniformément la charge.~» Il répartit uniformément les clés. La charge par clé est une tout autre affaire.}
\end{solution}

\begin{solution}{iq:iv3:system-design:10}
Le \Cref{lst:iv3:sd:states} présente une réponse possible. Une annulation qui croise la dernière exécution est normale~: l'exécution rend l'ordre final, et le rejet d'annulation qui arrive ensuite doit être accepté sans protestation. Une annulation non demandée émise par la plateforme fait passer l'ordre à l'état annulé, depuis n'importe quel état actif. Une exécution reçue après un accusé d'annulation, ou une exécution dépassant la quantité, déclenche une alarme.
\enlargethispage{-3\baselineskip}%
\omcode{code/interviews-3/06-system-design/python/iv_order_states.py}{1}{34}{La vie d'un ordre.}\label{lst:iv3:sd:states}

La vue de l'émetteur comporte des états en attente, car on ne sait rien tant que la plateforme n'a pas répondu. En pending new, l'ordre peut exister ou non. En pending cancel, il peut encore être exécuté, et le code continue de compter les exécutions à cet état. Il faut le traiter comme une exposition active~: tant que l'annulation n'est pas accusée, on peut encore négocier.

Un rejet d'annulation en pending cancel remet l'ordre à l'état ouvert. Cela signifie en général qu'il était trop tard, ou que l'identifiant était faux~; un système réel enregistre donc aussi la raison.

Les événements inattendus exigent une alarme, car ils signifient que notre image de la position ou de l'état de la plateforme peut être fausse. Une exécution après un accusé d'annulation signifie que l'un des deux se trompe, et la position peut être décalée. L'ignorer en silence laisse l'erreur grandir jusqu'à la fin de la journée.

Certaines plateformes rapportent autrement~: une exécution peut devancer l'accusé si les deux empruntent des canaux séparés, et on la traiterait alors comme un accusé implicite. Il faut donc lire la spécification de la plateforme et écrire la table en conséquence. Le test envoie des événements aléatoires à la table~: un état final reste final, et aucune sur-exécution ne passe.

\iqwrong{«~Un ordre a trois états~: ouvert, exécuté, annulé.~» Cela ne laisse aucune place à l'incertitude, et c'est pourtant là que se trouvent les pertes.}
\end{solution}

\begin{solution}{iq:iv3:system-design:11}
Le \Cref{lst:iv3:sd:risk} intègre les ordres ouverts dans le pire cas de chaque côté. Un contrôle sur la seule position courante est faux~: dix ordres ouverts, chacun assez petit pris isolément, peuvent tous être exécutés à la fois et dépasser largement la limite. Les contrôles se font en ligne, sur le seul chemin vers la plateforme, à coût constant.
\enlargethispage{-3\baselineskip}%
\omcode{code/interviews-3/06-system-design/python/iv_risk.py}{1}{33}{Contrôles avant transmission.}\label{lst:iv3:sd:risk}

Les deux côtés restent séparés, car ils ne se compensent qu'à l'exécution. Si l'on ne détient aucune position et que l'on a 1000 achats et 1000 ventes ouverts, l'un ou l'autre peut être exécuté seul~; le pire long est donc de 1000, et le pire short aussi. Compenser les achats par les ventes masquerait ce point. Les flux aléatoires du test vérifient l'assertion elle-même après chaque étape~: si tous les ordres ouverts étaient exécutés à cet instant, la position resterait dans la limite.

Ces contrôles doivent se trouver là où aucun ordre ne peut les contourner, entre l'ensemble des stratégies et la plateforme. Aux États-Unis, un courtier ayant un accès au marché doit disposer de tels contrôles sous son contrôle direct et exclusif. Si chaque stratégie garde les siens, un bogue dans une stratégie est aussi un bogue dans son contrôle.

Sur le chemin critique (hot path), chaque contrôle est une comparaison avec des valeurs déjà en mémoire, sans verrou, allocation ni réseau. On garde les champs en entiers, car la virgule flottante n'a pas sa place dans une limite. Son coût doit rester faible, bien qu'il s'applique à chaque ordre, et il faut le mesurer. Les prix de référence doivent être à jour, car un prix périmé fausse le collar~; il faut donc aussi le vérifier. Plus lent, le calcul du risque agrégé sur l'ensemble des stratégies tourne à côté de ce chemin et fixe ces limites.

\iqwrong{«~Vérifier la position après chaque exécution.~» À ce moment-là, l'ordre a déjà été négocié, et la limite est déjà franchie.}
\end{solution}

\begin{solution}{iq:iv3:system-design:12}
Un kill switch arrête d'abord les nouveaux ordres, puis annule ce qui est ouvert, puis prévient les personnes. Il peut être déclenché par une personne ou par une règle automatique. Il ne doit pas dépendre du processus qu'il arrête, puisqu'un processus défaillant est précisément le cas pour lequel il existe. Il évite de nouveaux dommages, sans défaire les transactions.

L'ordre des opérations compte. Il faut bloquer les nouveaux ordres en un point que tout ordre doit traverser, comme la passerelle ou la couche de risque~; sinon, une stratégie continue de remplacer ce que l'on annule. Il faut ensuite annuler tous les ordres ouverts, au moyen d'une annulation de masse (mass cancel) si la plateforme en offre une. La fermeture de la session peut suffire, sur une plateforme dotée d'une annulation à la déconnexion (cancel on disconnect). Il faut alors alerter les personnes, et laisser le système bloqué jusqu'à ce qu'une personne le redémarre.

Le kill switch doit s'exécuter sur son propre chemin, avec sa propre connexion, ou au moins dans un processus séparé. Un interrupteur placé dans la stratégie partage ses bogues. Les personnes doivent disposer d'un moyen simple de l'actionner, et du droit de le faire sans demander la permission.

Une ordonnance américaine de sanction décrit un broker-dealer qui a perdu plus de 460~millions de dollars en environ 45~minutes, à cause d'ordres qu'un ancien chemin de code continuait d'envoyer. Elle relève que la société ne disposait d'aucune procédure pour arrêter ce système en réponse à ses propres alertes. Il faut donc le répéter, car un interrupteur jamais testé risque de ne pas fonctionner le jour venu.

Déclencheurs automatiques utiles~: perte ou position au-delà d'une limite, cadence des ordres ou des messages au-delà de la normale, trop de rejets, perte des données de marché, et heartbeat qui ne vient plus. Le kill switch ne supprime pas la position déjà détenue, qui requiert toujours un plan de clôture, et il peut causer des dégâts à lui seul si l'annulation laisse une couverture inachevée. Il faut donc aussi trancher qui ferme quoi.

\iqwrong{«~La stratégie contrôle sa propre perte et s'arrête elle-même.~» C'est la stratégie qui peut être précisément ce qui est défaillant.}
\end{solution}

\begin{solution}{iq:iv3:system-design:13}
On met en tampon les mises à jour jusqu'à l'arrivée d'un snapshot, puis l'on n'applique que celles dont le numéro est postérieur à celui-ci, comme le fait le \cref{lst:iv3:sd:sync}. Si le snapshot est plus ancien que la plus ancienne mise à jour tamponnée, les numéros intermédiaires sont perdus~; il faut donc attendre un snapshot plus récent. La garantie~: chaque fois qu'un carnet existe, il est exact pour son numéro de séquence.
\omcode{code/interviews-3/06-system-design/python/iv_book_sync.py}{1}{35}{Snapshot et mises à jour.}\label{lst:iv3:sd:sync}

Chaque mise à jour fixe la taille d'un niveau, et une taille de 0 retire le niveau. Appliquer les mises à jour dans l'ordre à partir d'un point de départ correct donne donc un carnet correct. Une mise à jour de numéro inférieur ou égal à celui du snapshot y est déjà incluse et est ignorée~; c'est pourquoi le snapshot doit porter son numéro. Sans cela, on ne peut pas raccorder les deux flux.

Un trou, une fois synchronisé, signifie que le carnet n'est plus juste. Il ne faut pas continuer à l'afficher. On le retire, on prévient quiconque en dépend, et l'on repart du snapshot suivant. Un carnet faux en silence est pire que pas de carnet du tout, car une stratégie peut négocier contre des prix qui n'existent pas.

Le test exécute 300 flux aléatoires avec quelques pertes, fournit des snapshots légèrement en retard, et vérifie après chaque message que, si un carnet existe, il est égal au vrai carnet pour son numéro de séquence. Il couvre aussi un snapshot plus ancien que le tampon.

Deux points du listing demanderaient du travail en production. Le tampon n'a pas de limite tant que l'on n'est pas synchronisé, et la reprise après un trou peut prendre du temps. Un gestionnaire réel tenterait aussi d'abord une retransmission pour un petit trou, et ne prendrait un snapshot qu'en cas d'échec~; il garderait aussi les carnets de tous les instruments séparés, car un trou ne touche parfois qu'un seul d'entre eux.

\iqwrong{«~Appliquer le snapshot, puis toutes les mises à jour tamponnées.~» D'anciennes tailles absolues écraseraient des données plus récentes contenues dans le snapshot, et donneraient un niveau faux.}
\end{solution}

\begin{solution}{iq:iv3:system-design:14}
Le \Cref{lst:iv3:sd:conflate} ne garde qu'une valeur par symbole, de sorte que la mémoire est bornée par le nombre de symboles. Une mise à jour remplace la valeur et conserve la place du symbole dans la file, si bien qu'un symbole très actif ne peut pas repasser avant qu'un symbole calme ait eu son tour. Les transactions, les accusés d'ordre et tout autre flux d'événements ne doivent jamais être conflatés.
\omcode{code/interviews-3/06-system-design/python/iv_conflate.py}{1}{15}{Un tampon de conflation.}\label{lst:iv3:sd:conflate}

Il repose sur le fait qu'un dict Python conserve l'ordre d'insertion~: remplacer une valeur laisse la clé en place, tandis que la retirer puis la rajouter la renvoie en queue. Le test vérifie l'ordre de livraison par rapport à une simple liste de ceux qui sont devenus périmés en premier. Il vérifie aussi que le consommateur finit avec la valeur la plus récente de chaque symbole, quel que soit le rythme des messages.

La conflation sacrifie l'historique au profit de la fraîcheur. Un consommateur qui ne veut que le cours courant ne perd rien dont il ait besoin. Il préfère disposer du cours le plus récent plutôt que de traiter un arriéré de cours anciens~; c'est l'échec d'une simple file, où chaque cours affiché est aussi ancien que la file est longue.

Les données qui ne doivent pas être conflatées sont celles où chaque élément compte pour lui-même. Les transactions s'additionnent en volume, et les exécutions en une position. En remplacer une par une plus récente la fait simplement disparaître. Les mises à jour de carnet forment un cas particulier~: conflater jusqu'à l'état courant de chaque niveau est acceptable, puisque chacune est une taille absolue, mais abandonner les deltas incrémentaux ne l'est pas.

\iqwrong{«~Utiliser une file bornée, et abandonner la plus ancienne.~» Cela abandonne le symbole le plus ancien, même lorsque c'est sa seule mise à jour.}
\end{solution}

\begin{solution}{iq:iv3:system-design:15}
Le débit net de remplissage du tampon est de $5000 \times 100 = 0.5$ Mo/s, soit $16 \times 2^{20} / 500\,000 = 33.6$ secondes. Les choix sont alors de bloquer l'éditeur, d'abandonner des données, de conflater, ou de déconnecter le client. Conflater pour le tableau de bord. Déconnecter le système de risque pour qu'il se resynchronise, et alerter. Un moteur de trading ne devrait jamais être aussi lent, et doit échouer bruyamment.

Le blocage est le seul choix à écarter. Si l'éditeur attend un client lent, tous les clients sont retenus, si bien que le consommateur le plus lent fixe le rythme de tous. Chaque client doit donc avoir son propre tampon, et l'éditeur écrit sans bloquer.

Abandonner des données au sein d'un flux TCP casse le découpage ou la séquence sur lesquels le client s'appuie. Les abandons doivent donc se faire aux frontières de message, avec un numéro de séquence qui permette au client de voir le trou.

Pour un tableau de bord, une personne regarde quelques cours~; une valeur la plus récente, conflatée, est donc ce qu'il lui faut, sans perte pour personne. Un système de risque a besoin d'une image juste et complète~; une perte silencieuse est donc inacceptable. Une déconnexion accompagnée d'une alerte, puis une reconnexion et une resynchronisation à partir d'un snapshot, est une solution honnête. Un moteur de trading qui accuse 30~secondes de retard ne devrait plus négocier du tout~; le bon résultat est donc qu'il cesse de coter, et qu'une personne en soit informée.

Le tampon ne fait que retarder le problème. Si le client est lent en moyenne, un tampon plus grand ne fait qu'augmenter l'âge des données avant qu'il soit plein.

\iqwrong{«~Porter le tampon à 1~Gio.~» Cela ne retarde le problème que d'une demi-heure, et les données ont alors une demi-heure de retard.}
\end{solution}

\begin{solution}{iq:iv3:system-design:16}
Tout envoyer demande $10^6 \times 100 \times 50 \times 8 = 40$ Gbit/s, et ne suivre que 10 cours la ramène à 4 Gbit/s. À 20\,000 connexions par serveur, cela fait 50 serveurs, chacun à 800 Mbit/s au maximum. Les mises à jour leur parviennent depuis une couche interne de publication--abonnement. Lorsqu'un serveur tombe en panne, ses clients se reconnectent à un autre et reçoivent un snapshot frais.

La conception est en couches. Les cours entrent dans une petite couche d'ingestion, qui calcule une valeur courante par instrument et la publie sur un bus interne. Les serveurs de connexion s'abonnent, maintiennent les WebSocket et diffusent en éventail. Ils sont sans état, à part la liste d'abonnements de chaque client~; on peut donc en ajouter derrière un répartiteur de charge.

Le regroupement aide beaucoup. Plutôt qu'une trame par cours, on envoie à chaque client une trame par seconde avec ses 10 cours, car un humain ne lit pas plus vite. Un lot conflaté évite aussi d'envoyer des mises à jour que personne ne pourrait voir. Les tampons par client sont bornés, et un client lent reçoit la dernière valeur et non un arriéré.

À la reconnexion, les 20\,000 clients d'un serveur tombé reviennent tous d'un coup, en effet de troupeau. Il faut donc ajouter un délai aléatoire côté client, et servir un snapshot depuis un cache et non depuis la couche d'ingestion. Les clients d'un cours affiché supportent de petites pertes, contrairement à un système de trading, et la page doit signaler qu'un cours est périmé si les mises à jour cessent.

Un cache du snapshot courant, derrière une simple requête web, sert aussi de repli aux clients dont le réseau bloque les WebSocket.

\iqwrong{«~Envoyer chaque mise à jour à chaque client dès son arrivée.~» C'est la conception à 40 Gbit/s, et elle transporte surtout des octets que personne ne voit.}
\end{solution}

\begin{solution}{iq:iv3:system-design:17}
Ce n'est pas forcément un bogue. Le plus souvent, les deux journaux ne désignent pas le même instant~: leurs horloges diffèrent, les horodatages peuvent provenir de points différents sur chaque chemin, et un client peut avoir eu du retard. Seule une comparaison fondée sur le numéro de séquence de la plateforme, sur des paquets capturés de la même façon, permet de trancher.

Les causes possibles, de la plus probable à la moins probable. Les horloges sont décalées de plus que l'écart observé. Un horodatage est pris à la carte réseau, et l'autre après le passage du feed handler. Un client a été plus lent pendant une rafale, si bien que sa file contenait des données plus anciennes. Les deux suivent des canaux différents de la plateforme, par exemple un flux direct rapide et un flux consolidé plus lent. L'un a un trou qu'il n'a pas encore détecté, ou il conflate. Ou bien il existe un vrai bogue dans l'un des constructeurs de carnet (book builder).

Pour trancher, il faut aligner les deux journaux sur le numéro de séquence de la plateforme, et non sur l'heure murale. Si, au même numéro de séquence, les deux carnets concordent, il ne s'agissait que d'un décalage temporel, et l'on peut lire de combien l'un est en retard. S'ils divergent à ce point, l'un des carnets est faux, et l'on peut rejouer les paquets capturés de ce client dans son feed handler pour voir où.

Il faut donc des captures brutes de paquets, ou au moins le numéro de séquence et un horodatage de réception stocké avec chaque modification du carnet. Sans cela, on discutera des horloges indéfiniment.

\iqwrong{«~Même heure murale, donc ils doivent concorder.~» Deux horloges ne concordent jamais exactement, et un même instant ne désigne pas le même endroit dans le flux.}
\end{solution}

\begin{solution}{iq:iv3:system-design:18}
Elle affiche une part d'environ 0.1, puis \texttt{True True}. Dans notre exécution, la colonne des prix représente environ 10~\% du fichier compressé d'une journée. Les contrôles indiquent que le plan ne lit que la colonne \texttt{price}, et seulement un fichier sur quatre. Un format en colonnes permet à une requête de ne toucher que ce dont elle a besoin.

Un stockage en lignes devrait lire chaque champ de chaque ligne pour calculer la moyenne d'un seul. Dans Parquet, chaque bloc de colonne est stocké de façon contiguë, si bien qu'un lecteur ne prend que ceux qu'il veut, et le pied de fichier indique où ils se trouvent. Des valeurs de même nature se compressent aussi très bien ensemble. L'élagage des partitions fait un travail semblable sur les lignes, puisque les répertoires \texttt{date=} écartent des fichiers avant même de les ouvrir.

Ce format convient à l'analytique~: les balayages, les agrégats et quelques colonnes sur de nombreuses lignes. Il ne convient pas à la récupération de lignes entières isolées, ni à la mise à jour sur place, car les fichiers sont immuables et toute modification impose de les réécrire.

Pour l'organisation d'une année, on partitionne par date, puisque la plupart des requêtes désignent une plage de dates et que les jours anciens ne changent jamais. Dans le fichier d'une journée, on trie par symbole puis par heure, de sorte que les statistiques de chaque groupe de lignes permettent à une requête d'ignorer les groupes qui ne concernent pas son symbole. Il n'est pas nécessaire d'avoir un fichier par symbole. Pour des milliers de symboles, dont la plupart sont peu actifs, cela ferait un très grand nombre de minuscules fichiers, et ouvrir un fichier coûte bien plus que lire un bloc. Il faut aussi conserver les captures brutes, car elles constituent le seul enregistrement fidèle, et tout le reste peut être reconstruit à partir d'elles.

\iqwrong{«~Tout mettre dans une table relationnelle avec un index.~» Un index retrouve des lignes, alors que les requêtes d'un base de ticks lisent presque tout un intervalle de dates.}
\iqfollow{Et pour un symbole très actif~?}{Les statistiques des groupes de lignes s'en chargent, puisqu'un symbole très actif occupe simplement plusieurs groupes.}
\end{solution}

\begin{solution}{iq:iv3:system-design:19}
\Cref{lst:iv3:sd:delta} est un seul codec. 10\,000 horodatages occupent 20\,007 octets, contre 80\,000 à l'état brut~: 9 octets pour l'écart de la première valeur à zéro, puis 2 pour chacune des suivantes. Les heures voisines sont proches, si bien que les écarts sont petits, et un petit nombre demande peu d'octets. Des valeurs aléatoires occupent davantage, environ 9 octets chacune.
\omcode{code/interviews-3/06-system-design/python/iv_delta.py}{1}{26}{Écarts et entiers à longueur variable.}\label{lst:iv3:sd:delta}

Ici, un écart vaut 1000~ns. Le zigzag le double pour obtenir 2000, ce qui place le signe dans le bit de poids faible, et les petites valeurs négatives restent petites. 2000 demande 11 bits, soit 2 octets à raison de 7 bits par octet. Le codec est donc passé de 8 octets par valeur à 2 sur ces données, en tirant parti de leur structure, et le test confirme la taille exacte.

Avec des valeurs aléatoires, les écarts sont aussi grands que les valeurs elles-mêmes, le zigzag ajoute un bit, et le drapeau de 1 bit par octet coûte un huitième, si bien que la forme encodée dépasse 8 octets par valeur. Le test montre cette croissance. Aucun codec n'aide donc dans tous les cas, et celui-ci suppose des données triées.

Un magasin applique d'abord cette étape, car un compresseur général cherche des chaînes d'octets répétées, et une suite de nombres de 8 octets qui ne cessent de croître en contient peu. Une fois l'écart calculé, la même petite valeur se répète, et un compresseur général réussit alors très bien. Parquet dispose d'un encodage de cette famille, intégré pour les colonnes d'entiers.

Les horodatages réels de ticks sont moins réguliers que cela, si bien que 2 octets représente un cas haut de gamme pour la compression, et il faut mesurer sur ses propres données.

\iqwrong{«~Compresser le fichier avec gzip, et c'est suffisant.~» Il repère moins de structure qu'un codec qui sait que ce sont des heures triées, et sa relecture coûte davantage de CPU.}
\end{solution}

\begin{solution}{iq:iv3:system-design:20}
\Cref{lst:iv3:sd:bitemp} conserve chaque version avec l'heure à laquelle elle est devenue connue, et \texttt{as\_of} retient la plus récente connue à cette date. Un écrasement détruit ce que l'on savait, si bien qu'un backtest négocierait sur une correction qu'il ne pouvait pas avoir vue, et que le rapport du mois dernier ne pourrait plus être expliqué.
\omcode{code/interviews-3/06-system-design/python/iv_bitemporal.py}{1}{18}{Deux heures par ligne.}\label{lst:iv3:sd:bitemp}

Les deux heures répondent à des questions différentes. L'heure de l'événement indique quand quelque chose s'est produit, et l'heure de connaissance indique quand on a appris cette version de l'information. Un prix de lundi, corrigé le mardi, a une seule heure d'événement et deux heures de connaissance. Une question du type «~quel était le prix de lundi~?~» dépend donc du moment où on la pose.

Dans un backtest, la simulation de lundi après-midi doit utiliser ce qui était connu lundi après-midi. Si elle lit la valeur corrigée, l'anticipation est faible, mais elle est systématique, et elle flatte le résultat~: une stratégie évite de mauvaises cotations qu'elle aurait réellement traitées. Cela vaut à plus forte raison pour les résultats des sociétés et pour la composition des indices, qui sont révisés ou rétroalimentés.

Pour les rapports, les régulateurs ou un client peuvent demander ce que l'on croyait à l'époque et pourquoi on a agi ainsi. Avec des écrasements, personne ne peut répondre. Les requêtes as-of rendent aussi un résultat passé reproductible.

Le listing parcourt toutes les lignes, pour rester simple. Un vrai magasin trie par clé et par heure de connaissance, et ne lit que la dernière version à la date demandée ou avant. Le test compare cela au rejeu de toutes les corrections jusqu'à chaque instant, sur 300 tables aléatoires. Le stockage ne croît qu'avec les corrections, qui sont peu nombreuses.

\iqwrong{«~Garder une copie quotidienne de la table.~» Cela donne une résolution d'un jour, et coûte une copie complète chaque jour. Une correction en cours de journée resterait invisible.}
\end{solution}

\begin{solution}{iq:iv3:system-design:21}
Les prix nets se trient ainsi~: B 10.001, A 10.003, C 10.008. Le plan est donc 150 sur B, 100 sur A et 50 sur C, pour un total de 3000.85. Le \Cref{lst:iv3:sd:route} parcourt tous les niveaux par prix net. La réponse suppose que les tailles affichées seront toujours là quand les ordres arriveront, et c'est là le point faible.
\enlargethispage{-3\baselineskip}%
\omcode{code/interviews-3/06-system-design/python/iv_route.py}{1}{19}{Une répartition par prix net.}\label{lst:iv3:sd:route}

Le coût de chaque action est son prix plus les frais, ou moins la rétrocession. Comme les coûts s'additionnent sans interaction ici, prendre à chaque fois l'action la moins chère donne le meilleur total, et le test compare ce résultat à l'essai de toutes les répartitions entre trois plateformes.

Sur un marché réel, ce n'est qu'une première étape. D'autres voient aussi ces prix, si bien qu'un ordre passif visé peut avoir disparu à votre arrivée, surtout si vous êtes lent, ou si votre ordre adressé à une plateforme informe les autres de ce qui arrive. La taille affichée peut aussi ne représenter qu'une partie de ce qui est disponible. Un routeur réel envoie donc ses ordres de manière à ce qu'ils arrivent au même instant, en calant chaque plateforme sur sa latence. Il estime la probabilité d'exécution de chaque ordre, et pas seulement son prix. Il doit ensuite décider de la suite pour le reliquat~: le réacheminer, ou le laisser reposer dans le carnet.

Une rétrocession versée à celui qui prend une cotation attire aussi d'autres preneurs, si bien que cette cotation peut disparaître en premier, et un bon prix net assorti d'une faible chance d'exécution peut ne pas être une bonne affaire. Les règles comptent également~: les règles américaines exigent qu'un centre de négociation dispose de politiques interdisant de traverser une cotation protégée meilleure ailleurs. Le routeur ne doit jamais non plus envoyer au total plus que l'ordre parent, à moins d'en accepter le risque.

\iqwrong{«~Envoyer tout à la plateforme au meilleur prix.~» Ici, cette plateforme n'a que 150 à ce prix, et les frais modifient de toute façon le classement.}
\end{solution}

\begin{solution}{iq:iv3:system-design:22}
\Cref{lst:iv3:sd:recon} renvoie les trois écarts. \texttt{missing}~: la plateforme a une exécution que nous n'avons jamais vue. \texttt{extra}~: nous avons enregistré une exécution que la plateforme n'a pas. \texttt{differ}~: les deux l'ont, mais avec des détails différents. \texttt{missing} est le plus urgent, car votre position est fausse, et vos contrôles de risque en partent.
\omcode{code/interviews-3/06-system-design/python/iv_recon.py}{1}{8}{Exécutions comparées à la copie de la plateforme.}\label{lst:iv3:sd:recon}

Une exécution manquante représente une exposition réelle dont le système n'a pas connaissance, si bien que les limites et les couvertures sont fausses dès maintenant. Elle peut venir d'une session perdue, ou d'un bogue qui a fait tomber un message. Une exécution en trop signifie que votre position est fausse dans l'autre sens~: il peut s'agir d'un message de test qui a fuité, ou d'un doublon. Un écart de prix ou de quantité peut résulter d'une mauvaise lecture, d'une annulation d'exécution, ou d'une correction de la plateforme, et il touche votre P\&L avant de toucher votre risque.

Pourquoi pendant la journée~? Un écart découvert à la clôture a eu des heures pour grandir, et des stratégies ont négocié à partir de lui. La copie d'une plateforme arrive généralement par un flux séparé ou par un fichier, si bien qu'on peut lancer ce contrôle toutes les quelques minutes, ou en continu. Chaque écart doit déclencher une alerte avec un responsable, et la résolution doit être consignée.

Ce qui rend le contrôle exact, c'est la comparaison sur l'identifiant d'exécution. Comparer seulement les totaux réussirait si deux erreurs se compensent. Le calendrier rend toutefois un écart net difficile à établir~: une exécution proche de la limite de la fenêtre peut ne figurer que d'un seul côté pour le moment, si bien qu'un contrôle en direct doit revérifier un écart après un court délai avant de déclencher une alarme.

Le test supprime, modifie et ajoute des lignes aléatoires, et vérifie que chacune est signalée, et dans la bonne liste. Il couvre aussi le cas de deux côtés vides.

\iqwrong{«~Les positions concordent à la fin de la journée, donc tout va bien.~» Deux erreurs qui se compensent concordent aussi, et le risque intrajournalier était faux de toute façon.}
\end{solution}

\begin{solution}{iq:iv3:system-design:23}
Il faut enregistrer des écritures dont la somme est nulle, dans une seule transaction, comme dans le \cref{lst:iv3:sd:ledger}. Ainsi, l'argent ne peut ni se perdre ni se créer. Les nouvelles tentatives sont sûres grâce à la table \texttt{done}, et un plantage ne laisse rien à moitié fait. Un montant doit être positif, sinon un montant négatif retirerait de l'argent du compte cible et passerait le contrôle de solde. Le contrôle et les écritures vivent dans une même transaction, prise en écriture dès la première lecture, faute de quoi deux virements pourraient passer tous les deux le contrôle.
\omcode{code/interviews-3/06-system-design/python/iv_ledger.py}{1}{38}{Un grand livre sur SQLite.}\label{lst:iv3:sd:ledger}

Les écritures constituent l'historique~: un solde en découle, il peut donc toujours être vérifié, et un litige peut être remonté jusqu'aux lignes qui l'ont causé. Une colonne de solde stockée est un cache pour la rapidité, tenu dans la même transaction, et il doit être possible de la reconstruire à partir des écritures.

Le piège de la concurrence réside dans l'écart entre la lecture d'un solde et l'écriture. Avec deux virements de 80 prélevés sur un compte de 100, les deux pourraient lire 100 et passer tous deux. \texttt{BEGIN IMMEDIATE} prend le verrou d'écriture de SQLite avant la lecture, si bien que le second attend, puis voit 20. Sur une autre base, la même idée se traduit par un verrou de ligne avec \texttt{SELECT FOR UPDATE}, ou par un niveau d'isolation adapté. Ce n'est pas automatique, car beaucoup de bases utilisent par défaut un niveau qui laisse subsister cette course.

La ligne \texttt{done} est validée avec les écritures, si bien qu'une nouvelle tentative après un plantage trouve soit rien, et l'exécute, soit la première réponse. Les refus sont aussi enregistrés, de sorte qu'une requête a une réponse unique pour toujours~: une nouvelle tentative ultérieure d'un virement refusé reste refusée, ce que vérifie le test. Huit threads du test retirent chacun 30 d'un solde de 70, et exactement deux réussissent.

SQLite n'admet qu'un seul écrivain à la fois, il s'agit donc d'un modèle et non d'un dimensionnement. Les montants sont des entiers exprimés dans la plus petite unité, car la virgule flottante ne peut pas représenter exactement les centimes.

\iqwrong{«~Mettre à jour le solde sur place, dans une transaction.~» Cela fonctionne jusqu'au jour où quelqu'un demande pourquoi le solde est ce qu'il est.}
\end{solution}

\begin{solution}{iq:iv3:system-design:24}
Le 99\up{e}~centile d'une étape vaut 100. Celui du total vaut 120, et non 300. On ne peut donc pas additionner des centiles. La somme n'est juste que si les étapes sont lentes ensemble, auquel cas le centile du total vaut bien 300. Le \Cref{lst:iv3:sd:budget} calcule la distribution exacte, et le test la vérifie aussi par simulation.
\omcode{code/interviews-3/06-system-design/python/iv_budget.py}{1}{21}{La distribution d'une somme.}\label{lst:iv3:sd:budget}

Avec des étapes indépendantes, les totaux se répartissent ainsi~: 30 avec une probabilité de 94.1~\%, 120 avec 5.76~\%, 210 avec 0.12~\%, et 300 avec 0.0008~\%. Le 99\up{e}~centile tombe dans le deuxième groupe. Deux étapes lentes à la fois sont rares, si bien que la somme suppose une coïncidence qui ne se produit presque jamais. La médiane vaut 30 et la moyenne 35.4.

Cela joue dans les deux sens. Additionner les centiles des étapes surestime de 2.5 fois le cas indépendant. Mais une étape seule n'est lente que 2~\% du temps, alors que le total compte au moins une étape lente 5.9~\% du temps. La queue du total est donc atteinte bien plus souvent que celle de n'importe quelle étape. Une étape peut paraître correcte à son propre 95\up{e}~centile, et pourtant constituer ce dont est faite le 95\up{e}~centile de bout en bout.

Les moyennes s'additionnent, que les étapes soient indépendantes ou non, si bien qu'un budget de moyennes est sûr. Pour les queues, il faut mesurer le total, avec un horodatage à chaque étape sur un même message. Dans les systèmes réels, les étapes sont souvent lentes ensemble, par exemple à cause d'une cause commune comme une rafale, un vidage de cache ou une pause, et l'indépendance ne tient pas. Le cas corrélé du test donne 300, si bien que les chiffres réels se situent entre ces deux bornes.

\iqwrong{«~Le budget est de 50 par étape au 99\up{e}~centile, donc celui du total est de 150.~» Cela n'est vrai que si les trois sont lentes sur les mêmes messages.}
\end{solution}

\begin{solution}{iq:iv3:system-design:25}
Le temps moyen dans le système vaut $1 / (1 - \rho)$~: 2, 10 et 100 microsecondes. La simulation donne environ 1.99 et 9.74, et à 99~\% de charge 79, car 400\,000 échantillons ne suffisent pas à converger à ce niveau. Avec un service à durée fixe, l'attente est divisée par deux, si bien que le temps dans le système tombe à 1.5, 5.5 et environ 50. Il faut donc faire tourner un fil critique bien en dessous de sa capacité, et garder des durées de service régulières.
\omcode{code/interviews-3/06-system-design/python/iv_queue.py}{1}{14}{Un serveur, par simulation.}\label{lst:iv3:sd:queue}

La récurrence n'a besoin d'aucune liste d'événements. L'attente d'un client est celle du précédent, plus le service du précédent, moins l'écart entre les arrivées, et jamais en dessous de 0. Une seule ligne par client suffit donc, et cela constitue une bonne vérification de toute formule.

La hausse est brutale. Passer de 90~\% à 99~\% de charge représente 10~\% de travail en plus et 10 fois le délai. Une exécution plus longue est nécessaire, car la file à 99~\% met très longtemps à oublier son passé. Le résultat à 400\,000 échantillons est nettement inférieur à la formule, ce qui est en soi une leçon~: un test court d'un système chargé le flatte. Le test indique seulement que ce cas dépasse 50, et n'affirme pas 100.

La variabilité du service compte autant que la charge. L'attente est proportionnelle à $(1 + c^2)$, où $c$ est le coefficient de variation du service. Un service à durée fixe donne donc la moitié de celle d'un service exponentiel, et de rares services longs, comme un chemin lent occasionnel, aggravent beaucoup l'attente. Supprimer ces rares cas lents peut aider davantage que de rendre la moyenne plus rapide.

La moyenne masque aussi la queue de distribution. À charge donnée, le 99\up{e}~centile de l'attente est très au-dessus de la moyenne, et les objectifs sur la queue obligent en général à baisser la charge davantage que ne le font les objectifs sur la moyenne.

\iqwrong{«~Un système occupé à 90~\% laisse une marge de 10~\%.~» En termes de délai, la marge est déjà épuisée.}
\end{solution}

\begin{solution}{iq:iv3:system-design:26}
Un flush par ordre donne au plus 500 par seconde. Avec le group commit, chaque flush prend tout ce qui est arrivé depuis le lancement du précédent. À 10\,000 ordres par seconde, cela fait un lot de 20 et une latence moyenne d'environ 3~ms. À 300 par seconde, chaque lot ne contient qu'un ordre, et la latence se réduit au flush de 2~ms lui-même.
\enlargethispage{-3\baselineskip}%
\omcode{code/interviews-3/06-system-design/python/iv_group_commit.py}{1}{18}{Commit groupé.}\label{lst:iv3:sd:group}

Chaque ordre attend en moyenne la moitié d'un flush avant que le précédent se termine, puis son propre flush, soit $1 + 2 = 3$~ms. La taille du lot est le taux d'arrivée multiplié par la durée du flush. Un lot de 20 multiplie le débit par 20 pour un coût d'environ 1~ms. Ce qui séduit, c'est qu'il s'adapte~: sous faible charge, il n'y a aucune attente pour former un lot, comme dans le cas à 300 par seconde, et sous forte charge, le lot grossit de lui-même. Avec un flush par ordre à 10\,000 par seconde, la file croît sans limite, et dans le test, la latence moyenne presque double quand la durée de l'exécution double.

Ne pas acquitter tant que le flush couvrant cet ordre n'est pas terminé, car la promesse est que l'ordre se trouve sur le disque. Notons que 3~ms est très long pour un chemin rapide~: cette conception convient à un stockage plus lent et durable, et ce n'est pas ce qu'un moteur à faible latence place en ligne.

Une conception sensible à la latence peut remplacer le flush par une réplication~: un ordre est considéré comme sûr lorsqu'il est en mémoire sur une machine de secours. Cela coûte un aller-retour réseau au sein d'un même site, à la place du temps d'écriture sur disque. Cela change ce contre quoi l'on est protégé~: la panne d'une machine, et non la coupure de courant des deux. L'écriture sur disque a alors lieu ensuite, de manière asynchrone. Un stockage plus rapide modifierait aussi ces chiffres, il faut donc mesurer sur son propre système.

\iqwrong{«~Le batching ajoute de la latence, donc ne jamais grouper.~» Ici, c'est le batching qui permet au système de tenir du tout à ce débit.}
\end{solution}

\begin{solution}{iq:iv3:system-design:27}
La clé doit contenir la version exacte du code, une version exacte des données et tous les paramètres, sous une forme canonique, comme dans le \cref{lst:iv3:sd:runkey}. Les erreurs courantes sont les noms mobiles comme \texttt{latest}, les dictionnaires non triés et les paramètres oubliés. Une clé doit aussi prendre en compte les versions des bibliothèques et les graines. Ce qu'elle ne peut pas saisir, quelle que soit sa construction, ce sont les éléments qui ne sont pas du tout des entrées, comme les écarts de calcul en virgule flottante entre processeurs, ou l'ordonnancement des threads.
\omcode{code/interviews-3/06-system-design/python/iv_run_key.py}{1}{11}{Une clé de cache.}\label{lst:iv3:sd:runkey}

La règle est qu'une clé doit changer dès que le résultat pourrait changer. La version du code est donc un hash de commit issu d'un arbre propre, et non un nom de branche, car une branche bouge, et les modifications non validées sont invisibles. Le listing refuse le nom mobile \texttt{latest} pour les données. Une forme canonique empêche deux ensembles de paramètres égaux de produire des clés différentes, et \texttt{sort\_keys} s'en charge, ce que vérifie le test. Les flottants demandent de la prudence~: \texttt{0.1 + 0.2} et \texttt{0.3} sont des entrées différentes, donc un paramètre calculé peut rater le cache, ce qui est sûr, même si c'est coûteux.

Le sens dangereux est celui d'une clé qui ne change pas lorsque le résultat change. Un résultat périmé ressemble alors à un résultat réel, et personne ne s'en aperçoit. Il faut donc inclure les versions des bibliothèques, l'image de l'environnement et une graine pour tout ce qui est aléatoire, ou les intégrer à la version du code au moyen d'un environnement figé.

Certaines entrées ne peuvent pas entrer dans une clé, même si l'on peut s'en prémunir~: les écarts de calcul en virgule flottante entre types de processeurs, ou l'ordonnancement des threads dans un code non déterministe. Pour vérifier, il faut de temps en temps réexécuter un échantillon de résultats en cache et les comparer. Une divergence signifie que la clé oublie quelque chose.

\iqwrong{«~Hacher le nom de la fonction et ses arguments.~» La modification suivante du code renvoie l'ancien résultat, sans avertissement.}
\end{solution}

\begin{solution}{iq:iv3:system-design:28}
Par ordre de vraisemblance approximatif~: les données ont changé, par des corrections, des reconstitutions ou un nouveau rattachement des symboles~; l'univers ou les opérations sur titres ont changé~; un environnement ou une bibliothèque a changé~; une graine aléatoire n'a pas été fixée~; une valeur par défaut de configuration a bougé~; et le code n'était pas réellement identique. Il faut concevoir avec des versions de données figées, des environnements figés et un enregistrement durable de chaque exécution.

Les données sont la cause la plus fréquente, car les fournisseurs révisent l'historique, et une table nommée \texttt{prices} n'est plus aujourd'hui ce qu'elle était hier. Les entrées ont donc besoin de versions immuables, les nouvelles données arrivant sous forme de nouvelle version, et jamais en écrasant l'ancienne. Une exécution nomme alors exactement ce qu'elle a lu.

L'environnement vient en second~: le comportement par défaut d'une bibliothèque numérique ou d'une bibliothèque de dataframes peut changer d'une version à l'autre sans qu'aucune erreur ne soit levée. Chaque exécution doit donc utiliser une image figée ou un fichier de verrouillage, et le consigner. Pour le code, il faut un hash de commit issu d'un arbre propre, et les notebooks exécutés dans le désordre sont un piège fréquent.

Chaque exécution enregistre un relevé~: version du code, versions des données, environnement, paramètres, graine, type de machine et résultat. Une réexécution peut alors être comparée, et toute différence peut être tracée en modifiant une seule entrée à la fois. Avec une entrée figée, un résultat différent désigne une source de non-déterminisme restante, comme les threads ou les écarts de calcul en virgule flottante entre processeurs.

On peut donc reproduire un résultat tant que les entrées existent encore. Cela a un coût de stockage, et décider pendant combien de temps conserver les anciennes versions de données est un véritable choix de politique. Une explication suffit souvent~: «~le ratio de Sharpe a baissé de 0.2 depuis que le fournisseur a retiré ces 40 noms, et voici le changement.~»

\iqwrong{«~Les réexécutions diffèrent parce que les marchés sont aléatoires.~» Un backtest est un calcul sur des données stockées, donc une différence a une cause.}
\end{solution}

\begin{solution}{iq:iv3:system-design:29}
Avec un seul worker, la pire latence de la plateforme saine atteint 433~s, et elle continuerait de croître dans une exécution plus longue. Avec un worker par plateforme, elle vaut 10~ms, comme si la plateforme bloquée n'existait pas. Les ordres en cours sur la plateforme bloquée sont dans un état inconnu, et on les compte donc comme une exposition vivante tant que le contraire n'est pas établi.
\omcode{code/interviews-3/06-system-design/python/iv_bulkhead.py}{1}{13}{Un worker partagé, ou un par plateforme.}\label{lst:iv3:sd:bulkhead}

C'est l'isolation, parfois appelée cloisonnement (bulkhead)~: une défaillance dans une partie ne doit pas se propager à travers une ressource partagée. Une file est la ressource partagée habituelle, avec un pool de threads, un verrou et un tampon de socket. Des workers par plateforme donnent à la plateforme saine du test exactement les mêmes chiffres que lorsqu'elle tourne seule. Le modèle ne borne pas la file de la plateforme bloquée, alors que dans un système réel une borne doit exister, faute de quoi la mémoire se remplit.

Un routeur a donc aussi besoin d'une limite de temps sur la réponse de chaque plateforme, et d'un disjoncteur. Après un nombre donné de délais dépassés ou de heartbeats manqués, il faut marquer la plateforme comme indisponible, cesser de lui router des ordres, et rejeter ou réacheminer les nouveaux ordres. Le réacheminement a toutefois un revers. Les ordres ouverts sur la plateforme bloquée peuvent être exécutés à tout moment, si bien qu'en envoyant la même quantité ailleurs, on risque d'être exécuté deux fois. Il faut donc compter ces ordres dans sa limite de position comme s'ils pouvaient être exécutés.

Pour lever l'état inconnu, il faut tenter une annulation si le chemin le permet, demander un état après une reconnexion, et lire le drop copy s'il est disponible. Les réglages d'annulation à la déconnexion aident là où ils sont proposés. Rien n'est tranché tant que la plateforme ne l'a pas confirmé, et en dernier ressort, une personne devra peut-être téléphoner à la plateforme.

\iqwrong{«~La plateforme ne répond plus, donc ces ordres sont morts.~» Ils peuvent très bien être vivants chez la plateforme, et être négociés.}
\end{solution}

\begin{solution}{iq:iv3:system-design:30}
Le \Cref{lst:iv3:sd:journal} écrit une longueur, un CRC puis les octets, et la reprise s'arrête au premier enregistrement qui ne se vérifie pas. La longueur indique où se termine l'enregistrement, et le CRC indique si ses octets sont bien ceux qui ont été écrits. Un rejeu complet de 5 milliards d'événements à un million par seconde prend 83~minutes, on conserve donc des instantanés.
\omcode{code/interviews-3/06-system-design/python/iv_journal.py}{1}{23}{Un journal.}\label{lst:iv3:sd:journal}

Un plantage peut laisser un dernier enregistrement tronqué~: seule une partie a atteint le disque, ou le fichier s'est allongé au-delà d'octets qui n'ont jamais été écrits. La longueur seule ne révélerait pas des données parasites qui tombent par hasard à la bonne taille, et c'est à cela que sert le CRC. Une suite de zéros demande de la prudence~: le CRC de rien vaut 0, donc avec une graine nulle, huit octets nuls passeraient pour un enregistrement vide. Une graine non nulle règle ce problème. La reprise tronque alors le fichier à la longueur renvoyée avant d'ajouter des données, sans quoi les enregistrements suivants se placeraient après les mauvais octets et seraient perdus à la reprise suivante.

Le test coupe le fichier à chaque octet, et attend exactement les enregistrements complets situés avant la coupure. Il ajoute aussi des zéros et inverse des bits au hasard, et attend que la reprise s'arrête au niveau de l'enregistrement endommagé ou avant, et jamais après. Une corruption ne peut donc pas produire un faux enregistrement, sauf si le CRC lui-même entre en collision, ce qu'un CRC sur 32~bits rend improbable mais possible à ce volume. Une erreur au milieu du fichier signifie une perte de données, et non un plantage, si bien qu'un système réel devrait la signaler par une alarme, et non tronquer le fichier en silence.

Pour la reprise, il faut enregistrer l'état à intervalles réguliers, marqué par la position dans le journal, puis charger le dernier instantané et rejouer ce qui suit. Avec un instantané toutes les 5~minutes, le rejeu porte au plus sur 5~minutes d'événements. Si les 5 milliards arrivent régulièrement sur une session de 6.5~heures, cela fait environ 64 millions, soit une minute de travail. Même si les événements arrivaient aussi vite que le rejeu s'exécute, cela ferait 300 millions, soit 5~minutes. Il faut prendre les instantanés sur une machine de secours, afin que le serveur principal ne se mette pas en pause. Ce format ne dit rien de la durabilité~: celle-ci exige un flush avant chaque acquittement.

\iqwrong{«~Chaque écriture est atomique, donc un enregistrement est soit présent, soit absent.~» Une écriture dans un fichier peut être interrompue en plein milieu par un plantage.}
\end{solution}

\begin{solution}{iq:iv3:system-design:31}
Avec un seul thread, il n'y a pas de verrous, l'ordre des événements est celui de l'exécution, et les mêmes entrées donnent les mêmes sorties. À 2~microsecondes par ordre, un thread atteint au plus 500\,000 par seconde, on partitionne donc par instrument. Un ordre qui touche deux partitions perd son atomicité. Le déterminisme est ce qui rend la réplication simple.

Les verrous coûteraient plus que le travail, car un appariement touche plusieurs structures partagées, dans un ordre fixé. Chaque ordre doit de toute façon attendre son tour dans le carnet, il n'y a donc rien à exécuter en parallèle au sein d'un même carnet. Un thread unique n'a pas non plus de lignes de cache disputées, et un chemin très prévisible. Une conception publiée fait tourner sa logique métier de cette façon, et garde tout son état en mémoire, avec un journal durable des entrées.

Partitionner par instrument ou par groupe d'instruments, un thread chacun. Chaque carnet est indépendant, donc cela passe à l'échelle avec le nombre de carnets, et des instruments liés, comme un contrat à terme et ses options, peuvent rester sur un même thread s'ils doivent interagir. Il reste alors le carnet le plus sollicité comme limite, qu'il faut repousser en accélérant le code. Répartir un seul carnet sur plusieurs threads coûte plus qu'il ne rapporte.

Un ordre de spread qui traverse des partitions demande une coordination~: une réservation en deux étapes dans chaque partition, un propriétaire unique pour les carnets du spread, ou un carnet de spread distinct qui négocie en legging sur les instruments simples. Chacune de ces solutions coûte de la latence ou de la complexité, ce qui est une raison de placer ces instruments ensemble.

Pour la réplication, un principal et une machine de secours n'ont besoin que de la même séquence d'entrées. Avec un moteur déterministe, ils se retrouvent alors dans le même état, sans mémoire partagée ni validation en deux phases. Leurs sorties peuvent être comparées pour contrôle.

\iqwrong{«~Utiliser davantage de threads pour aller plus vite.~» À l'intérieur d'un carnet, l'ordre est justement l'essentiel, donc les threads passeraient surtout à s'attendre les uns les autres.}
\end{solution}

\begin{solution}{iq:iv3:system-design:32}
Appliquer puis enregistrer surcompte~: un plantage entre les deux fait qu'un redémarrage applique à nouveau cette entrée. Enregistrer puis appliquer sous-compte, car un plantage entre les deux fait sauter une entrée. Une étape unique et atomique est exactement ce qu'il faut. Le \cref{lst:iv3:sd:consumer} modélise les trois, et le test de plantage aléatoire détecte l'erreur des deux séquences incorrectes.
\omcode{code/interviews-3/06-system-design/python/iv_consumer.py}{1}{34}{Trois façons de consommer un journal.}\label{lst:iv3:sd:consumer}

Le test utilise 2000 journaux et des scénarios de plantage aléatoires. L'approche atomique donne toujours la vraie somme. Appliquer d'abord ne descend jamais en dessous, et dépasse parfois. Enregistrer d'abord ne dépasse jamais, et descend parfois en dessous. Il vérifie que chaque erreur se produit réellement, car un modèle incapable de montrer une erreur ne prouverait rien.

Une livraison exactement une fois ne peut pas être garantie à travers un réseau, car un émetteur ne peut jamais distinguer un message perdu d'une réponse perdue. Ce que l'on peut obtenir, c'est un effet exactement une fois, grâce à une livraison au moins une fois combinée à un effet idempotent ou atomique avec la position du consommateur. Ici, il faut conserver la position du consommateur dans la même transaction que l'état~: une ligne de base de données à côté du résultat.

Si l'état se trouve sur un système sans transaction partagée, il faut rendre l'application elle-même idempotente. On attache la position dans le journal ou un identifiant unique à chaque modification, et l'autre système refuse celle qu'il a déjà vue. Une API de virement bancaire avec une clé de requête en est un exemple. S'il n'offre ni l'un ni l'autre, il ne reste qu'à choisir laquelle des deux erreurs on peut accepter, puis à la détecter par un rapprochement. Pour un flux d'exécutions, on ne choisirait jamais de sous-compter en silence. Pour des métriques, surcompter un peu peut convenir.

Le modèle ne comporte qu'un seul consommateur, avec un état durable. Il ne couvre pas un second consommateur qui croit que le premier est tombé en panne, et c'est là qu'intervient la question du fencing.

\iqwrong{«~Mon courtier de messages offre l'exactly-once.~» Cela couvre en général son propre journal et ses propres transactions, et s'arrête à la base de données.}
\end{solution}

\begin{solution}{iq:iv3:system-design:33}
Il affiche \texttt{False 250}, puis \texttt{True 90}. Avec un fencing, le stockage refuse l'écriture de l'ancien principal. Sans lui, la valeur périmée 90 écrase 250. L'époque est un numéro qui augmente à chaque changement de principal, et la ressource écrite doit la contrôler, et non l'écrivain, qui ignore qu'il est périmé.
Un délai d'attente ne peut pas distinguer un processus mort d'un processus lent. Une longue pause, une partition du réseau ou une machine virtuelle bloquée ressemblent toutes à une mort, et le processus reprend ensuite sans aucune notion du temps écoulé. Il peut donc exister, à un instant donné, deux principaux, et aucun réglage du heartbeat ne l'élimine. Un délai plus court multiplie les fausses alertes.

La solution a deux volets. Une source convenue distribue les époques, comme un service de consensus, de sorte qu'il n'en existe qu'une par numéro. Et chaque écriture porte l'époque, le stockage retenant la plus haute qu'il a vue. Le contrôle de l'écrivain périmé a ainsi lieu au point d'effet, de manière atomique avec l'écriture, et la conviction de l'ancien principal n'a plus d'importance. Un algorithme de consensus appelle ce numéro un terme, et le décrit comme une horloge logique qui permet aux serveurs de repérer les leaders périmés.

Le cas plus difficile est un effet situé hors du stockage. Une plateforme ne vérifie pas l'époque, si bien qu'un ordre de l'ancien principal passe. Il faut donc une autre barrière, plus proche de la plateforme. Il peut s'agir d'une passerelle qui vérifie l'époque, et par laquelle tout ordre doit passer, ou de faire révoquer l'ancienne session par la bascule (failover), sur une plateforme qui n'accepte qu'une seule session par identifiant. Sinon, une annulation à la déconnexion peut retirer les ordres ouverts, mais seulement une fois que la plateforme a vu la déconnexion. Il faut s'attendre à un certain recouvrement, puis procéder à un rapprochement.

Le modèle du listing s'exécute dans un seul processus, il montre donc la règle et non une défaillance.

\iqwrong{«~Fixer un bail de 10 secondes~: l'ancien principal s'arrêtera de lui-même.~» Un processus suspendu pendant plus longtemps ignore que son bail a expiré, et ne le découvre qu'après avoir écrit.}
\end{solution}

\begin{solution}{iq:iv3:system-design:34}
Attendre tous les accusés de réception prend 40 ms, et une majorité prend 1 ms. Une majorité de 3 survit à une panne, alors qu'attendre tout le monde s'arrête à la moindre panne. Si le site à 1 ms tombe, la validation prend 40 ms. Deux majorités doivent se recouvrir, afin que la majorité d'un nouveau leader contienne toujours quelqu'un qui possède chaque écriture validée.
\omcode{code/interviews-3/06-system-design/python/iv_quorum.py}{1}{10}{Temps de validation et majorités.}\label{lst:iv3:sd:quorum}

Le temps de validation est une statistique d'ordre~: le $k$-ième plus petit aller-retour. Une majorité de 3 vaut 2, l'écrivain compris, donc le site à 40 ms ne ralentit pas une validation, alors qu'il finit tout de même par conserver une copie. Cela signifie aussi que le site éloigné peut être en retard à tout moment. La perte du site proche ne laisse que le site éloigné pour compléter une majorité, si bien que toutes les validations passent à 40 ms. C'est le point à anticiper, car il multiplie par 40 la vitesse de votre système précisément au moment où les choses vont déjà mal.

Une majorité de $n$ tolère $(n-1)/2$ pannes, arrondi à l'entier inférieur, si bien qu'une grappe de cinq en tolère deux. Le test vérifie que deux majorités quelconques ont un membre en commun, pour $n$ allant jusqu'à 7, en les énumérant toutes.

Avec moins d'une majorité, disons un seul accusé de réception, deux côtés d'une partition réseau peuvent valider chacun de leur côté, chacun croyant être le seul. Après la reconnexion, les deux historiques divergent, et une écriture validée est perdue. La réplication asynchrone fait délibérément le même compromis~: des validations plus rapides, avec une perte possible des dernières écritures lors de la bascule. Cela peut convenir à certaines données, si l'on sait réconcilier, mais c'est alors une décision, et non un accident.

Un chemin de négociation accepte souvent l'asynchronisme sur les longues distances, avec des copies synchrones seulement à proximité, car personne ne peut négocier avec 40 ms de latence sur chaque ordre. En cas de sinistre, l'état du site éloigné reste donc en retard, et il faut noter de combien.

\iqwrong{«~Trois copies signifient l'absence de perte de données.~» Seulement si la validation les attend. Des copies encore en chemin ne sont pas des copies.}
\end{solution}

\begin{solution}{iq:iv3:system-design:35}
Le gaspillage vaut environ $C/T + T/(2M)$, minimal en $T^* = \sqrt{2CM}$. Avec $C = 60$ s et $M = 86\,400$ s, on obtient $T^* = 3220$ s, soit environ 54 minutes, et un gaspillage de 3.7~\%. Une longue simulation, portant sur $10^8$ secondes de travail, donne 3.8~\% à cet endroit, et environ 8~\% au quart de cet intervalle comme à quatre fois celui-ci. Une exécution courte est plus bruitée. Inutile de s'en soucier pour un travail bien plus court que $M$.
\omcode{code/interviews-3/06-system-design/python/iv_checkpoint.py}{1}{18}{Points de contrôle et pannes, simulés.\protect\hfil\protect\break}\label{lst:iv3:sd:checkpoint}

Le premier terme est la part du point de contrôle dans chaque intervalle. Pour le second, une panne fait perdre en moyenne la moitié d'un intervalle, et les pannes surviennent au taux $1/M$, si bien que le temps perdu par unité de temps vaut $T / (2M)$. En annulant la dérivée, on obtient $C/T^2 = 1/(2M)$, d'où la formule. Au meilleur point, les deux termes sont égaux, ce qui constitue une vérification commode. Le minimum est plat~: un facteur 2 dans un sens ou dans l'autre ne coûte qu'environ un quart de gaspillage supplémentaire. Il n'est donc pas nécessaire de connaître $M$ exactement, et c'est pourquoi une valeur approximative suffit.

Aux extrêmes~: lorsque $M$ croît sans limite, $T^*$ croît et le gaspillage tend vers 0. Lorsque le coût $C$ tend vers 0, on crée des points de contrôle en permanence. Si $T^*$ n'est pas nettement plus petit que $M$, l'approximation échoue, car les pannes surviennent alors au sein d'un intervalle, et la formule simple sous-estime le gaspillage. La simulation n'a pas une telle limite, et c'est pourquoi nous l'exécutons~: sur notre travail, la formule et la simulation concordent à 15~\% près.

Inutile de s'en préoccuper lorsqu'un travail est court par rapport à $M$ et facile à relancer. Un travail de 10 minutes sur cette machine échoue dans environ 0.7~\% des cas, et le relancer coûte moins que n'importe quel schéma de points de contrôle. Découper le travail en petits morceaux indépendants lorsque c'est possible~: cela remplit le rôle d'un point de contrôle, avec moins de mécanique. Le taux de panne de l'ensemble du parc compte aussi, car un travail sur 100 machines échoue 100 fois plus souvent.

\iqwrong{«~Créer un point de contrôle aussi souvent que possible, par sécurité.~» Chacun coûte du temps, et au quart de l'intervalle optimal, dans cet exemple, le gaspillage vaut environ 8~\%.}
\end{solution}

\begin{solution}{iq:iv3:system-design:36}
Un ordre envoyé il y a quelques minutes ne doit pas repartir, car son prix peut être désormais loin du marché, et son émetteur peut y avoir renoncé. Un message de comblement de trou le remplace~: le numéro est conservé, et il ne porte rien. Un numéro ne peut pas simplement être sauté, car le récepteur verrait un trou et réclamerait à nouveau. Il faut ensuite interroger le statut de chaque ordre non résolu.
\omcode{code/interviews-3/06-system-design/python/iv_session.py}{1}{36}{Une session qui renvoie les messages.}\label{lst:iv3:sd:session}

La numérotation des sessions garantit que rien ne passe inaperçu. Si un numéro manque, le récepteur s'arrête et demande un renvoi. Chaque numéro doit donc être comptabilisé, même pour un message que l'on décide de ne pas envoyer. Un protocole très répandu offre cette fonctionnalité~: une réinitialisation de séquence en mode comblement de trou répond à une demande de renvoi, et fait avancer le récepteur au-delà des messages qui ne seront pas envoyés. Le listing conserve les messages de statut et envoie un comblement de trou pour les ordres plus anciens que \texttt{max\_age}.

Choisir \texttt{max\_age} est une vraie décision. Un ordre renvoyé peut dater de quelques secondes et rester valable, ou au contraire être dangereux, selon le marché. Beaucoup de systèmes vont plus loin, et ne renvoient jamais automatiquement un ordre après une reconnexion~: ils les comblent tous, et laissent la stratégie décider.

Le test exécute 500 sessions aléatoires avec pertes. Il vérifie qu'aucun ordre périmé n'est accepté, que chaque numéro est traité, et que les messages acceptés sont dans l'ordre, sans répétition. Il couvre aussi le cas d'un message reçu deux fois, que le récepteur ignore.

Que devient donc un ordre qui n'a pas été renvoyé~? On ne peut pas supposer qu'il n'est jamais arrivé, car seule sa réponse a pu être perdue. L'ordre est donc dans un état inconnu, et il faut interroger la plateforme ou lire le drop copy~; ce n'est qu'alors que l'on connaît sa position. D'ici là, il faut le compter comme potentiellement actif.

\iqwrong{«~Repartir simplement à un nouveau numéro de séquence à chaque reconnexion.~» Alors un rapport d'exécution perdu disparaît, et la position devient fausse sans bruit.}
\end{solution}

\begin{solution}{iq:iv3:system-design:37}
Déployer d'abord sur un seul serveur, dans un mode qui ne modifie aucun comportement, puis élargir par paliers, avec une personne qui surveille et un retour arrière testé. Les indicateurs de fonctionnalité démarrent désactivés, et s'activent séparément. Ne jamais réutiliser un ancien indicateur, car l'ancien code peut lui donner son ancienne signification. Deux versions côte à côte doivent pouvoir lire les messages et l'état l'une de l'autre.

Avant l'ouverture, il faut vérifier que chaque serveur possède bien la version attendue, en interrogeant chacun d'eux, et non en faisant confiance à l'outil qui a procédé au déploiement. Comparer aussi la configuration, et envoyer des ordres de test par un symbole de test ou par l'environnement de test d'une plateforme. Une décision de sanction américaine décrit un cas~: le nouveau code avait été installé sur sept serveurs sur huit, sans contrôle d'une seconde personne. Le huitième conservait un vieux code mort, et la nouvelle version réutilisa l'indicateur de l'ancien code. Ce seul serveur émit des millions d'ordres non voulus, et la société perdit plus de 460~millions de dollars en environ 45 minutes. Les leçons à en tirer~: supprimer le code mort, vérifier tous les serveurs, et ne pas réutiliser les indicateurs.

Activer d'abord une nouvelle fonctionnalité pour une seule stratégie, un seul symbole ou une petite taille, et laisser l'indicateur séparé de la mise en production du code. Cela garde «~déployer~» et «~activer~» comme deux décisions distinctes, si bien qu'éteindre une fonctionnalité est rapide et ne nécessite aucune mise en production.

Pour revenir en arrière, il faut disposer d'un chemin connu et répété à blanc. Cela est plus difficile qu'il n'y paraît si une nouvelle version a écrit un état que l'ancienne ne sait pas lire. Les changements d'état et de messages se font donc par étapes~: d'abord, faire en sorte que les lecteurs acceptent les deux formes, puis modifier les écrivains, et enfin retirer plus tard l'ancienne forme.

Limiter les dégâts grâce à des limites d'ordres et à un kill switch fonctionnant indépendamment de la mise en production, et éviter les mises en production aux heures où le marché est le plus actif.

\iqwrong{«~Il a passé tous les tests, donc le déployer partout à l'ouverture.~» Dans ce cas, la défaillance ne se trouvait pas dans le code testé, mais dans ce qui tournait réellement.}
\end{solution}

\begin{solution}{iq:iv3:system-design:38}
Les passerelles contrôlent les ordres de chaque membre, un séquenceur donne un ordre total, et cette séquence est journalisée et répliquée. Un moteur d'appariement mono-thread l'applique ensuite, puis publie les résultats. Même journal, même état~: le cœur est donné par le \cref{lst:iv3:sd:exchange}. 100\,000 ordres par seconde représentent 10 microsecondes chacun, si bien qu'un thread par partition peut y suffire, à condition que le journal suive.
\omcode{code/interviews-3/06-system-design/python/iv_exchange.py}{1}{32}{Le cœur déterministe d'un moteur d'appariement.\protect\hfil\protect\break}\label{lst:iv3:sd:exchange}

Dans l'ordre~: la passerelle authentifie, analyse, et applique les contrôles de risque au niveau du membre ainsi que les limites de débit, chacun sur son propre état de session, et elle ajoute l'identifiant d'ordre client du membre. Le séquenceur attribue le numéro et une heure, et écrit le journal, avec une copie vers un secours en veille installé sur le même site, avant tout accusé de réception. Le moteur d'appariement lit le journal et émet des événements, chacun portant le numéro de son entrée. Les éditeurs envoient les accusés de réception d'ordres aux membres, et les données de marché à tous.

L'équité est la règle selon laquelle l'ordre du séquenceur est l'ordre de priorité. La priorité prix-temps est donc équitable si les chemins vers le séquenceur sont égaux, et il faut publier ce que l'on fait de ces chemins~: mêmes longueurs de câble, même traitement, et personne ne bénéficiant d'une route privée. Sinon, la priorité temporelle récompense celui dont le câble est le plus court.

Pour un secours exact, chaque entrée du moteur d'appariement doit provenir du journal, y compris l'heure. Le jumeau du test lit une horloge à l'intérieur du moteur, et un rejeu du même journal donne alors une empreinte différente. Les nombres aléatoires, l'ordre de parcours des tables de hachage et les calculs flottants qui diffèrent d'une machine à l'autre relèvent du même danger. Il faut comparer les empreintes entre le principal et le secours, comme vérification continue.

Pour se rétablir~: repérer le dernier enregistrement du journal, charger le dernier instantané, puis rejouer. Le secours dispose déjà de tout cela, il peut donc simplement prendre le relais, sous une époque plus élevée pour écarter l'ancien. Rien de perdu~? Le journal contient tout ce qui a été accusé de réception, puisque rien n'est accusé avant d'être durable. Comparer les numéros de séquence avec ceux des passerelles et du flux publié. Puis republier à partir de la séquence que les membres ont vue en dernier. Informer les membres de l'état du marché, et rouvrir par une enchère si les règles l'exigent.

Le test vérifie l'absence de carnet croisé, l'équilibre des quantités, et une empreinte inchangée après une bascule à un point aléatoire, sur 100 journaux aléatoires. Le système réel, avec son réseau et ses temps de réponse, n'a pas été exécuté.

\iqwrong{«~Horodater à la passerelle, et trier selon cet horodatage.~» Les horloges diffèrent d'une passerelle à l'autre, si bien que cet ordre n'en est pas un, et une arrivée tardive obligerait à réécrire l'historique.}
\end{solution}

\begin{solution}{iq:iv3:system-design:39}
Le débit de pointe vaut $2 \times 10^6 \times 64 \times 8 = 1.02$ Gb/s en multicast, une seule fois, ou 307 Gb/s en unicast vers les 300 destinataires. On utilise donc le multicast en interne, et l'unicast à la périphérie. Une session de 6.5 heures au débit moyen représente environ 750 Go bruts, et ce même débit sur 24 heures donnerait 2.76 To. Les moteurs de négociation reçoivent directement le multicast, le risque prend un flux en conflation, la recherche capture tout, et les tableaux de bord prennent un petit flux en conflation.

L'architecture~: une paire de feed handlers par ligne de plateforme, sur le même site que la source, chacun normalisant vers un seul format et une seule séquence internes. Ils publient chaque groupe d'instruments sur un groupe multicast, de sorte qu'un consommateur ne rejoint que ce dont il a besoin. Les groupes les plus chargés reçoivent leurs propres cœurs et matériels, et la répartition doit reposer sur des débits mesurés.

Les pertes se jouent sur deux niveaux. Du côté de la plateforme, les numéros de séquence et l'arbitrage de lignes détectent une perte, et la reprise de la plateforme la répare. En interne, chaque flux possède sa propre séquence, et un service de rejeu conserve le flux récent pour la retransmission, avec un service d'instantanés pour les consommateurs qui arrivent tard ou accusent trop de retard. Un taux de perte de 1 sur 10\,000 au pic représente 200 messages par seconde, si bien qu'un système de retransmission a besoin de sa propre capacité, et d'une limite, au cas où un problème réseau généralisé ferait demander tout le monde en même temps.

Pour la bascule, il faut faire tourner en direct les deux feed handlers de chaque paire, côte à côte. Les consommateurs arbitrent entre eux, ou bien un principal désigné publie et un moniteur effectue la bascule. Quand une ligne de plateforme tombe, seules les données de cette plateforme deviennent périmées~: il faut en informer les consommateurs par un message de statut explicite, instrument par instrument, car une ligne silencieuse ressemble à un marché silencieux. Les heartbeats sur chaque flux permettent de les distinguer.

Mesurer les débits par ligne, les trous de séquence, le nombre de reprises et la latence de chaque étape, avec des horodatages à la carte réseau et à la publication, ainsi que le retard de chaque consommateur. Capturer les paquets bruts à l'entrée, car c'est l'enregistrement à rejouer après un bogue. Ce dimensionnement repose sur les données de l'énoncé~: aucun système n'a été exécuté, et la latence réelle doit se mesurer sur le matériel considéré.

\iqwrong{«~Un seul grand bus de messages pour tout le monde.~» Un moteur de négociation et un tableau de bord n'ont pas les mêmes besoins, et sur un chemin partagé, le plus lent impose ses conditions à tous.}
\end{solution}

\begin{solution}{iq:iv3:system-design:40}
Données~: 5 Gio $\times$ 252 jours $\times$ 10 ans $\times$ 25 régions font 315\,000 Gio, soit environ 308 Tio, avant les versions et les copies. Temps~: 500 $\times$ 3 $\times$ 8 minutes font 12\,000 minutes-cœur, soit environ 3.1 heures sur 64 cœurs pour une étude. Couches~: données brutes, données nettoyées point-in-time, variables, simulation, résultats. Et un bon résultat doit avoir résisté au hors-échantillon et à l'examen des coûts.

La couche brute conserve ce qui est arrivé, de manière immuable, partitionnée par date et par région, en colonnes. Les données nettoyées sont dérivées par des traitements versionnés, et chaque publication reçoit une version, si bien que les corrections arrivent comme de nouvelles versions, avec leurs dates de connaissance. Les données de référence, comme les changements de symbole, les opérations sur titres et l'appartenance aux indices, sont stockées par date as-of. C'est là qu'intervient le biais de survie, si bien que la recherche les lit au moyen d'une seule bibliothèque, et non à la main.

La simulation tourne sur un ordonnanceur, avec un travail par jour de stratégie, ce qui rend les nouvelles tentatives peu coûteuses, puisqu'un morceau en échec se relance simplement. Une étude isolée prend 3 heures lorsque le parc est libre, et 50 chercheurs ne tiendront pas tous en même temps~: il faut donc des quotas et des priorités. Si chacun lance une étude de ce type par jour, la demande vaut $50 \times 200$ heures-cœur, soit 10\,000, alors que le parc en fournit 1536 en 24 heures. Cela explique qu'il faille soit un parc plus grand, soit de la capacité spot, soit un cache de résultats, indexé sur la version exacte du code, la version des données et les paramètres. Le cache rapporte souvent en premier, puisque la plupart des relances changent peu de chose.

Les coûts font l'objet d'une bibliothèque commune, avec un modèle versionné, car un backtest qui fixe ses propres coûts n'est pas comparable. Les hypothèses sur les exécutions et sur la file d'attente doivent être énoncées, et le simulateur doit être testé contre les exécutions réelles.

Pour prévenir un accident, il faut enregistrer automatiquement chaque exécution, y compris celles qui ont échoué, car le nombre de tentatives compte lorsqu'on juge un ratio de Sharpe. Réserver les données récentes, et en contrôler l'accès. Exiger le même chemin de code pour le réel et pour la recherche, et comparer les résultats réels au backtest après le lancement. Ce dimensionnement repose sur les hypothèses de l'énoncé.

\iqwrong{«~Donner à chaque chercheur une grosse machine et une copie des données.~» Les copies divergent, les résultats ne peuvent pas être comparés, et personne ne sait quel chiffre faire confiance.}
\end{solution}
